% SPDX-License-Identifier: CC-BY-4.0
% Copyright (C) 2026 Fabian Januszewski, Abhijit (the Langlandscape team)
\documentclass[11pt,a4paper]{article}
\usepackage[T1]{fontenc}
\usepackage{lmodern}
\usepackage[margin=2.2cm]{geometry}
\usepackage{amsmath,amssymb,amsthm}
\usepackage{enumitem}
\usepackage{booktabs}
\usepackage{array}
\usepackage{xcolor}
\usepackage[most]{tcolorbox}
\usepackage{tikz}
\usetikzlibrary{arrows.meta,positioning}
\usepackage[colorlinks=true,linkcolor=blue!50!black,urlcolor=blue!50!black]{hyperref}

\setlist{itemsep=2pt,topsep=3pt}
\setlength{\parindent}{0pt}
\setlength{\parskip}{4pt}

% notation
\newcommand{\Q}{\mathbb{Q}}
\newcommand{\Z}{\mathbb{Z}}
\newcommand{\Zp}{\mathbb{Z}_p}
\newcommand{\Qq}{\mathbb{Q}_q}
\newcommand{\cO}{\mathcal{O}}
\newcommand{\T}{\mathbb{T}}
\newcommand{\cT}{\mathcal{T}}
\newcommand{\cM}{\mathcal{M}}
\newcommand{\cX}{\mathcal{X}}
\newcommand{\Bhat}{\widehat{B}}
\newcommand{\XB}[1]{X_B(#1)}
\newcommand{\GL}{\mathrm{GL}}
\newcommand{\SL}{\mathrm{SL}}
\newcommand{\Nrd}{\mathrm{Nrd}}
\newcommand{\Sym}{\mathrm{Sym}}
\newcommand{\BG}{\mathrm{BG}}
\newcommand{\cc}{\mathfrak{c}}
\newcommand{\tq}{\mathfrak{t}}
\newcommand{\m}{\mathfrak{m}}
\newcommand{\ord}{\mathrm{ord}}
\newcommand{\Frob}{\mathrm{Frob}}
\newcommand{\Phicomb}{\Phi^{\mathrm{comb}}}
\newcommand{\coker}{\operatorname{coker}}
\newcommand{\HypX}{\mathrm{Hyp.X}}

% status tags
\newcommand{\st}[1]{\,{\small\fcolorbox{gray!60}{gray!8}{\textsc{#1}}}}
\newcommand{\proved}{\st{proved}}
\newcommand{\cited}{\st{cited}}
\newcommand{\pmc}{\st{proved mod.\ cited}}
\newcommand{\toprove}{\st{to prove}}
\newcommand{\expected}{\st{expected}}
\newcommand{\conj}{\st{conjectural}}

\newcommand{\role}[1]{\textbf{Role.} #1}

\title{\textbf{Langlandscape: the six-stage plan}\\[2pt]
\large A human-readable summary of \texttt{plan/staged/}}
\author{}
\date{15 November 2026 --- final state: all six stages completed}

\begin{document}
\maketitle
\vspace{-1.5em}

%======================================================================
\section*{0.\quad The big picture}

The project rebuilds the \textbf{definite Jacquet--Langlands correspondence over $\Q$} as an explicit, integral
and \emph{simplicial} construction. Throughout, $B=B_D$ is the definite quaternion algebra over $\Q$ ramified exactly
at $\infty$ and at the primes dividing a squarefree $D$ (with an odd number of prime factors). The basic idea:
\begin{itemize}
  \item On the quaternion side, automorphic forms are just functions on a \textbf{finite set}, the class set
        $\XB{U}=B^\times\backslash\Bhat^\times/U$.
  \item Choose an auxiliary prime $q\nmid D$. The norm-one group $\cO_U[1/q]^1$ acts on the \textbf{Bruhat--Tits tree}
        of $\SL_2(\Qq)$, and the quotient is a finite graph built out of class sets.
  \item The homology of these graphs plays the role of the character group of a modular or Shimura curve with bad
        reduction at $q$. On these graphs the classical arguments of Ribet (level raising) and Ihara (Ihara's lemma)
        become combinatorics.
  \item Everything is then taken up the $p$-power tower and made $\Lambda$-adic (Hida theory).
\end{itemize}

The plan has three layers:
\begin{tcolorbox}[colback=gray!5,colframe=gray!50,boxrule=0.5pt,arc=1pt]
\textbf{(a) Combinatorial core.} Tree quotients, the definite Ribet sequence, the definite Ihara lemma. This needs only
strong approximation and graph theory.

\textbf{(b) One geometric hypothesis, $\HypX$.} It identifies the graph objects with $\GL_2$ objects (character groups
of Jacobians). It is used \emph{only} to interpret and to produce maps.

\textbf{(c) Arithmetic output, independent of $\HypX$.} For $D=q$, a $\Lambda$-adic \textbf{defect formula}
\[
  (h_F)=\cc_q(F)=\tq_q(F),\qquad h_F=\eta_F/\eta_F^B,
\]
comparing the congruence ideals on the $\GL_2$ side and the quaternion side with the old/new congruence ideal $\cc_q$
and the Tamagawa ideal $\tq_q$ (generated by $\rho_F(I_q)-1$). This extends the fixed-weight results of Kim--Ota (and
Pollack--Weston in weight~2) along a whole Hida family.
\end{tcolorbox}

Classical JL, Chenevier's $p$-adic JL and the $\Lambda$-adic Hecke-algebra JL (Longo--Vigni, Hsieh) are \textbf{cited,
not reproved}. Whether the last arrow of the construction (from the Deligne--Rapoport graph of the modular curve at $q$
to its Jacobian) can itself be made simplicial is left as an open question.

\begin{center}
% Status colours (update when a stage's status changes):
%   done = completed (merged on main), prog = in progress, todo = not started.
\definecolor{stDone}{HTML}{CDEBC5}\definecolor{stDoneD}{HTML}{2E7D32}
\definecolor{stProg}{HTML}{FFE6A8}\definecolor{stProgD}{HTML}{B07800}
\definecolor{stTodo}{HTML}{EEEEEE}\definecolor{stTodoD}{HTML}{8A8A8A}
\tikzset{stbox/.style={draw,rounded corners=2pt,inner sep=4pt,font=\small,align=center,line width=0.8pt},
  done/.style={stbox,fill=stDone,draw=stDoneD},
  prog/.style={stbox,fill=stProg,draw=stProgD},
  todo/.style={stbox,fill=stTodo,draw=stTodoD,dashed}}
\begin{tikzpicture}[node distance=0.55cm and 0.55cm, >={Stealth}]
  \node[done] (s1) {01\\foundations};
  \node[done,right=of s1] (s2) {02\\tree quotients};
  \node[done,right=of s2] (s3) {03\\Ribet--Ihara};
  \node[done,right=of s3] (s4) {04\\$\Lambda$-adic};
  \node[done,above right=0.05cm and 0.7cm of s4] (s5) {05\\comparison $\HypX$};
  \node[done,below right=0.05cm and 0.7cm of s4] (s6) {06\\integral defect};
  \draw[->] (s1) -- (s2); \draw[->] (s2) -- (s3); \draw[->] (s3) -- (s4);
  \draw[->] (s4) -- (s5); \draw[->] (s4) -- (s6);
\end{tikzpicture}

\smallskip
{\small
\tikz\node[done,inner sep=2pt]{\phantom{x}}; completed\quad
\tikz\node[prog,inner sep=2pt]{\phantom{x}}; in progress\quad
\tikz\node[todo,inner sep=2pt]{\phantom{x}}; not started\quad
(status as of 15 Nov 2026, 04:05 GMT+1)}
\end{center}

{\small\textbf{Status tags} used below: \proved{} (in the plan), \cited{} (from the literature), \pmc,
\toprove, \expected, \conj. A \emph{hypothesis} is an explicit assumption; in Lean it is a named \texttt{Prop},
never an axiom.}

%======================================================================
\section{Stage 1 --- Quaternionic foundations}

\role{A conventions-and-bookkeeping layer with no deep new mathematics. It fixes every object, normalisation and sign on
the definite side once, so that the later stages cannot silently disagree.}

\subsection*{Objects}
\begin{itemize}
  \item A maximal order $\cO\subset B$ with splittings $\cO_\ell\cong M_2(\Z_\ell)$ for $\ell\nmid D$.
  \item Levels $U_0(N)$, $U_1(N)$ (lower row $\equiv(0,1)$), the $p$-tower $U_s=U_1(Mp^s)$ with class sets $S_s$, and
        $U_0(q)U$ (Iwahori at $q$).
  \item For $x\in\Bhat^\times$: the stabiliser $\Gamma_x=B^\times\cap xUx^{-1}$ (finite, of order dividing $24$) and
        the weight $w_x=|\Gamma_x/\pm1|$. In lattice terms, $x$ is a right ideal and $\Gamma_x$ is the unit group of its
        left order.
\end{itemize}

\subsection*{Results}
\begin{itemize}
  \item \textbf{Stabiliser lemma} \proved. For $p\ge5$: $\Nrd(\Gamma_x)=1$; at level $U_1(p)$ every $\Gamma_x$ is
        trivial and $\Delta_s=(1+p\Zp)/(1+p^s\Zp)$ acts freely up the tower; $p\nmid|\Gamma_x|$.
  \item \textbf{Mass formula.} Eichler's $\sum_x 1/|\Gamma_x|=\frac1{24}\prod_{\ell\mid D}(\ell-1)$ \cited, and its
        level change by the index $[U:U']$. With weights $w_x$ the level change holds iff $-1\in U'$ or $-1\notin U$;
        otherwise an extra factor $2$ appears (the \textbf{factor-2 hypothesis} used throughout). Consequence:
        \[|S_s|=\tfrac{\varphi(D)}{24}\,\psi(Mp^s)\,\varphi(Mp^s)\qquad(\text{e.g. }D=3,\ p=5:\ 2,\,50,\,1250).\]
  \item \textbf{Brandt modules} $\cM(U;V)$ with $V_k=\Sym^{k-2}$, the homological dual $\Z[X]$, the pairing
        $\langle[x],[y]\rangle=w_x\delta_{xy}$ and the norm map $\nu:[x]\mapsto w_x\mathbf 1_x$.
        \emph{Subtlety:} $\nu$ is Hecke-equivariant only up to adjoints ($T\mapsto T^*$), and it turns pushforward into
        trace.
  \item \textbf{Hecke operators} $T_\ell,S_\ell,U_\ell,\langle d\rangle,\Pi_\ell$ ($\ell\mid D$), $U_p$.
        \emph{Main decision:} $U_p=\sum_{j=0}^{p-1}\left(\begin{smallmatrix}p&j\\0&1\end{smallmatrix}\right)$, the
        normalisation that commutes with pushforward up the $p$-tower \proved. The seed's literal formula
        $[U\,\mathrm{diag}(1,p)\,U]$ is the \emph{other} operator, the one that commutes with trace.
\end{itemize}

\subsection*{Convention flags passed on}
\begin{itemize}
  \item The adelic diamond $\langle d\rangle$ is the classical $\langle d\rangle^{-1}$.
  \item Under JL, $U_\ell\mapsto+\Pi_\ell$ for $\ell\mid D$. Checked on $D=11$ against the elliptic curve 11a; this
        contradicts the sign printed in Hsieh.
\end{itemize}

\subsection*{Cited targets}
Classical JL (Brandt form), Chenevier's $p$-adic JL, and the $\Lambda$-adic Hecke JL
$\T^{\ord,D\text{-new}}\cong\T^{\ord}_B$.

%======================================================================
\section{Stage 2 --- Tree quotients}

\role{Turn the finite class sets into a graph. Fix $q\nmid Dp$. The group $\Gamma_U=\cO_U[1/q]^1$ (norm-one elements,
integral away from $q$) acts on the Bruhat--Tits tree $\cT_q$, preserving types and without inversions.}

\subsection*{Results}
\begin{itemize}
  \item \textbf{Quotient graph} (\emph{proved given connectivity}). $\Gamma_U\backslash\cT_q$ is the
        \textbf{bipartite Brandt graph} $\BG(U,q)$: vertices are two copies of $\XB{U}$, edges are $\XB{U_0(q)U}$,
        endpoint maps $s_0[y]=[y]$ and $s_1[y]=[y\,\mathrm{diag}(1,q)]$, and stabilisers are the $\Gamma_x$. Hecke
        operators away from $q$ commute with $s_0,s_1$.
  \item \textbf{Deep inputs} \cited: strong approximation (SA) for $B^1$ at $\{\infty,q\}$, and the Hasse--Schilling
        norm theorem (used silently in the seed).
  \item \textbf{Connectivity criterion} (\emph{proved, new}). For one level $U$, SA at that level is equivalent to
        \textbf{$\BG(U,q)$ being connected} --- a finite check that HACS or Lean can do directly.
  \item \textbf{Two-term complex} \proved. Since $p\nmid|\Gamma_x|$, $H^*(\Gamma_U,V)$ is computed by
        \[0\to\cM(U;V)^2\to\cM(U_0(q)U;V)\to0.\]
        Harmonic cochains, the kernel of $(s_{0*},s_{1*})$, equal $H_1$ and are taken as the \textbf{definition of the
        $q$-new part}.
  \item \textbf{Free case.} At level $U_1(p^s)$, $\Gamma_U$ is free of rank $(q-1)|\XB{U}|+1$ and everything is plain
        graph homology.
  \item \textbf{Indefinite stand-in.} $\cX_{Dq}(U):=H_1(\BG,\Z)$ plays the role of the character group of the Shimura
        curve of discriminant $Dq$.
\end{itemize}

\subsection*{Caveats}
\begin{itemize}
  \item Over $\Z$, $H_1(\Gamma_U,\Z)$ has $2$- and $3$-torsion from the vertex groups (e.g.\ $D=2$, $q=3$ gives
        $\Z/3$ although the graph has $H_1=0$). Hence the graph's $H_1$ is used, not group homology.
  \item $H_0(\Gamma_U,V)$ needs genuine density, which no single-level check gives.
  \item SA is \textbf{not} available in Lean (the seed was wrong about this); it is a registered hypothesis.
\end{itemize}

%======================================================================
\section{Stage 3 --- Ribet and Ihara on the definite side}

\role{Prove the definite analogues of Ribet's level-raising exact sequence and of Ihara's lemma, as statements about
weighted graphs.}

\begin{itemize}
  \item \textbf{Weighted graphs.} For a graph with edge weights $w_e$: the lattice $X(\mathcal G)$, a weighted pairing,
        and the combinatorial component group $\Phicomb$. Two-vertex formula
        $|\Phi|=\prod_e w_e\cdot\sum_e 1/w_e$ \proved.
  \item \textbf{Mazur check} \pmc: $\Phicomb$ of the Deligne--Rapoport graph at $q$ is
        $\Z/\mathrm{num}\bigl(\tfrac{q-1}{12}\bigr)$, recovering the component group of $J_0(q)$.
  \item \textbf{$\Phicomb$ is Eisenstein} (\emph{proved for all $q$}): $T_\ell-(1+\ell)$ kills it (an elementary
        argument found in the audit).
  \item \textbf{Definite Ribet sequence} \pmc, obtained by collapsing vertices:
        \[0\to H_1\to\Z[\XB{U'}]^0\to\bigl(\Z[\XB{U}]^0\bigr)^2\to0.\]
  \item \textbf{Definite Ihara lemma.} Parts (i), (ii) \proved. Part (iii) --- $T_\ell$ acts on the cokernel as
        $(1+\ell)$, i.e.\ the cokernel is Eisenstein --- \toprove.
  \item \textbf{Level-raising operator.} $\det(\alpha\alpha^\dagger)=(q+1)^2-T_q^2$ \toprove{} --- the familiar
        level-raising congruence condition.
  \item \textbf{$\GL_2$ side.} Amalgam and Mayer--Vietoris, with $\det\Psi^*\Psi=(q+1)^2-q^{2-k}T_q^2$. The $\GL_2$
        Ihara lemma in Betti form is only \expected.
\end{itemize}

%======================================================================
\section{Stage 4 --- The $\Lambda$-adic passage}

\role{Run stages 2 and 3 up the $p$-tower $S_s$ and apply Hida's ordinary projector $e=\lim U_p^{n!}$ to obtain
statements over $\Lambda=\cO_E[[1+p\Zp]]$.}

\begin{itemize}
  \item \textbf{$e$ kills the junk} \proved: the degree term, the vertex term, $\Phicomb\otimes\Zp$ and $H_0$ all
        vanish after $e$. The ordinary pairing between $e$ and its adjoint $e^*$ is perfect (with $e\ne e^*$ in
        general).
  \item \textbf{Freeness and control} \pmc: the ordinary Brandt module $\cM_B^{\ord}$ is $\Lambda$-free and specialises
        correctly at each weight (elementary Hida-type argument; also in Hsieh and Longo--Vigni).
  \item \textbf{Weight points.} Because of the diamond inversion from stage 1, the seed's weight points must be inverted:
        $P_k^{\mathrm{ad}}=\ker\bigl([u]\mapsto u^{2-k}\bigr)$.
  \item \textbf{$\Lambda$-adic definite Ribet sequence}, exactness \proved:
        \[0\to\cX^{\ord}_{Dq}\to\cM^{\ord}_B(U_0(q)U)\to\cM^{\ord}_B(U)^2\to0.\]
        It needs connectivity at \emph{every} level $s$, which no finite check provides (so SA is genuinely used).
  \item \textbf{From Hsieh} \cited, under a residual condition: the localised Hecke algebra is Gorenstein, the Brandt
        module is free of rank one over it, and the quaternionic congruence ideal $\eta_F^B$ is a self-pairing.
\end{itemize}

%======================================================================
\section{Stage 5 --- The comparison hypothesis $\HypX$}

\role{Isolate, in one precise hypothesis, everything that compares the graph objects with the geometry of modular curves,
and derive its consequences.}

\subsection*{The hypothesis $\HypX(M;q)$}
For every $s$, an isomorphism between the character group at $q$ of the Jacobian $J_s$ and the degree-zero divisors on
the class set,
\[\varphi_s:\ X_q(J_s)\xrightarrow{\ \sim\ }\Z[S_s]^0,\]
such that
\begin{enumerate}[label=(X\arabic*)]
  \item it is Hecke-equivariant, with $U_q\mapsto+\Pi_q$;
  \item the monodromy pairing goes to the weighted pairing $w_e\delta$;
  \item it is natural in $s$;
  \item it respects the Galois filtration, with $\Frob_q=q\,U_q$.
\end{enumerate}
At $\Gamma_0$-level the ingredients are \cited{} (Deligne--Rapoport, Ribet 1990, SGA\,7). At $\Gamma_1(Mp^s)$-level
with $U_p$ there is essentially no reference: \expected. In Lean it is a named \texttt{Prop}, never an axiom.

\subsection*{Results}
\begin{itemize}
  \item \textbf{Sign} $\varepsilon=+1$ fixed \pmc, against Hsieh's printed $-$.
  \item The \textbf{$\Lambda$-adic version} of $\HypX$ follows formally; the Ohta twist is only sketched.
  \item A non-canonical isomorphism version holds at good maximal ideals \pmc.
  \item \textbf{Sub/quotient comparison} \proved: the $q$-new sub- and quotient lattices differ by $\T/(\mathfrak a+\mathfrak b)$.
  \item Frobenius conventions fixed once (Ribet--Stein, Picard functoriality): $\Frob_q=q\,U_q$.
\end{itemize}

%======================================================================
\section{Stage 6 --- The integral defect}

\role{Prove the arithmetic output (c): for $D=q$, the ratio of congruence ideals equals the old/new congruence ideal
equals the Tamagawa ideal, along the $q$-new Hida branch. This does \textbf{not} use $\HypX$.}

\textbf{Status: completed} (merged on \texttt{main} at \texttt{aed93c37}; final gate at Lean tree \texttt{4405815}:
PASS, PrintAxioms $836/836$ standard axioms, from scratch on TITAN and on CONCEPTD).

\begin{itemize}
  \item \textbf{Commutative algebra} \proved: Lemmas S, G, D and Theorem~E (i)--(iv),
        \[\eta_F=\cc_q\cdot\eta_F^B\]
        under Gorenstein hypotheses. Without the second Gorenstein hypothesis it holds at height-one primes, and $\cc_q$
        is principal iff that ring is Gorenstein. In Lean the layer \texttt{CommAlg/} is \emph{formalized} (standard
        axioms only, $0$ \texttt{sorry}), including the DVR half of Theorem~E(iii).
  \item \textbf{Galois half}: $\tq_q\subseteq\cc_q$, via Chebotarev and Burnside. The output of the Galois steps 1--3 is
        the registered hypothesis \texttt{Hyp.FreeRankTwo}, which is \emph{equivalent} to the conclusion
        $\tq_q\le\cc_q$ given \texttt{ResidualBurnside} (\texttt{Defect.freeRankTwo\_iff\_le}). \texttt{ResidualBurnside}
        is conditionally discharged from absolute irreducibility.
  \item \textbf{Squeeze.} At a Kim--Ota point, level lowering mod $\varpi^t$ gives the other inclusion, hence the
        \textbf{defect theorem}, \pmc{} (Hida control, Hsieh, Kim--Ota) and registered hypotheses. In Lean,
        \texttt{Defect.defect\_formula}, \texttt{cor\_fixed} and \texttt{cor\_allweights} are \emph{formalized modulo
        the registered hypotheses, including the Galois half (\texttt{Hyp.FreeRankTwo})}; this is not an unconditional
        theorem. The all-weights corollary needs $c\ne0$ at both points (not assumed).
  \item \textbf{The case $q\equiv-1$.} $M=1$: \pmc. $M>1$: F1 and F1$'$ (F1$'$ under our reading R4) \pmc, i.e.\ for
        squarefree $M$ with $\bar\rho$ unramified at all $\ell\mid M$. \textbf{$N(\bar\rho)>1$ is open}: the
        counterexample B2 refutes the step-(4) proof (and only that step).
  \item \textbf{Corrections to the seed:} the old quotient is $\T(M)_\m$ via $U_q\mapsto\alpha_q$ (the seed had twice
        the rank); one hypothesis needs ``no $p$-new forms''; $\cO_P$ is \emph{not} claimed to be a DVR.
  \item \textbf{Betti incarnation.} $\coker(\Psi^*\Psi)_\m\cong(\T(M)_\m/\delta)^2$ computes $\cc_q$ in rank one
        (HACS computation: see below).
  \item \textbf{Computed checks:} T4 Hecke side lengths $1,1,1,1,2$ at $k=12,28,44,60,124$, as predicted; definite
        Gram side positive control $(13,8,37)$ with $v_p=1$; the C4(i) regression is partial ($131/319$ cases run).
  \item \textbf{Several primes} ($D$ with $\ge3$ prime factors): $(h_F)=\prod_i\tq_{q_i}$ \conj.
\end{itemize}

%======================================================================
\section*{HACS (computations)}

All numbers in this section are \emph{computed}, with HACS v0.2.1
(DOI \href{https://doi.org/10.5281/zenodo.24231446}{10.5281/zenodo.24231446}) and hacs-data v0.2.0
(DOI \href{https://doi.org/10.5281/zenodo.24147012}{10.5281/zenodo.24147012}). \textbf{Lane A} (a Python bridge
driving HACS) and \textbf{Lane B} (exact integral congruence modules and a native quaternion core inside HACS) are both
closed and merged on \texttt{main}.
\begin{itemize}
  \item \textbf{Mazur on the definite side} ($k=2$, $B_{q,\infty}$). Through HACS's integral engine the Eisenstein
        congruence module equals $\Z/\mathrm{num}\bigl(\tfrac{q-1}{12}\bigr)$ \emph{exactly} ($2$- and $3$-parts
        included) on all $14$ levels of Lane~A's records ($B_{11}$ and every prime $D\le43$). The read of the native
        census confirms it for every prime $q\le600$ at $\ell\ge5$ ($109$ levels), and exactly for the first $80$ primes,
        $q\le409$ (run capped).
  \item \textbf{Beyond Mazur, $k>2$} ($32$ integral $\Sym^n$ cells). The only nontrivial lattice congruence module is at
        $D=13$, $k=8$: $|C|=37$. On the $\GL_2$ side, $2^3\cdot37$ is the only resultant of the $T_2$ factors that
        contains $37$. \emph{Caveat:} the lattice at $k>2$ is not canonical.
  \item \textbf{Two independent implementations} (HACS-native C++ and Python + cypari2) give the same $69/69$ orbit Smith
        lists. There is \textbf{no read-side speed claim}: both take milliseconds.
  \item \textbf{Build side} (class set + $T_2$ for $B_{997,\infty}$, $h=83$, single-threaded, same node): Python
        $64.7$\,s, native $0.02$\,s; ``${\sim}10^3\times$'' is the honest line.
  \item \textbf{Native census:} all $404$ primes $q\le2777$ generated in $19.6$ min; Eichler's class number formula holds
        for all $1229$ primes $D<10^4$. The full $404$-cell HACS read was not done.
  \item \textbf{Betti side of stage 6} (\texttt{T4\_betti}): $\coker(\Psi^*\Psi)_\m=(\Z/17)^2=(\T(1)_\m/\delta)^2$ at
        $k=12,28,44,60$; $k=124$ is not reachable. Rows with $k\ge\ell$ are tagged \emph{expected} (outside the
        amalgam's proven range $\ell>k$); the cross-check against S06-C3 is pending.
\end{itemize}
Validation was targeted only; one full HACS run (\texttt{-{}-all}, strict golden, digests) is due after the merge.

%======================================================================
\section*{Audits and open points}

Each stage of the plan had a Muse audit (no mathematical errors; mainly citation fixes), and a global Fable ``seam''
audit across all stages returned \textbf{pass with required fixes} (7 major, 13 minor items). In the implementation,
every stage had Fable audits of its Lean code and of its chapter, plus cross-account second readers. For stage~6: Lean
part~1 PASS WITH REMARKS $0/0/9$, part~2 $0/2/11$; chapter part~1 $0/2/12$, part~2 FAIL for the $M>1$ layer as
first written, then PASS WITH REMARKS $0/2/4$ on the restricted text, part~3 $0/0/7$. Open at the end:

\begin{center}
\small
\begin{tabular}{>{\raggedright\arraybackslash}p{0.68\textwidth}l}
\toprule
\textbf{Open point} & \textbf{Where}\\
\midrule
$q\equiv-1$ with $N(\bar\rho)>1$ (B2 refutes step (4) only); F1 is numerically untested & Stage 6, open\\
\texttt{Hyp.FreeRankTwo} (Galois half) and the other cited \texttt{Hyp.*}: no Lean term discharges them & Stage 6, registered\\
C4(i) regression partial ($131/319$); Eisenstein claim needs $\cO_v$ a DVR & Stage 6, computed\\
\texttt{prop:s05-Xlambda} (c2): six missing steps, no proof given & Stage 5, open\\
$\HypX$ (\texttt{Hyp.X.char}, \texttt{.iso}, \texttt{.gal}); the tower version needs \texttt{Hyp.MordLambdaFree} & Stage 5, registered\\
Definite Ihara lemma (iii) (\texttt{Hyp.IharaH0} in Lean) & Stage 3, registered\\
Kim--Ota Gorenstein claim vs.\ [BKM] at trivial primes (avoided by $q\not\equiv1$) & Stage 6\\
Several primes: $(h_F)=\prod_i\tq_{q_i}$ & Stage 6, conjectural\\
Simplicial version of the last arrow (graph $\leftrightarrow$ Jacobian) & open question\\
SA and Hasse--Schilling in Lean & registered hypotheses\\
Sign with Hsieh (FD2): \emph{decided} (Fabian: $U_q\mapsto+\Pi_q$ accepted; [Hs] prints $JL(U_q)=+U_q$, so Thm~4.5 as printed is vacuous); open only: contact the author & post-Mathathon\\
Citations not checkable on this node (B\"ockle 2007, Dimitrov 2009, EPW06, Edixhoven 1992, Wiles 1988) & unverified\\
Lean minors R9--R11; \texttt{linter.overlappingInstances} suppressed on 3 headlines & Lean, post-Mathathon\\
HACS: full validation run; canonical lattices for $k>2$; T4\_betti vs.\ S06-C3; full census read & HACS, post-Mathathon\\
\bottomrule
\end{tabular}
\end{center}

\section*{Lean status (final, on \texttt{main})}

The Lean layers of all six stages are merged on \texttt{main} (toolchain \texttt{v4.35.0}), with $0$ \texttt{sorry} of our own (the
only \texttt{sorry} is the one in the imported trace-formula project) and every reported declaration on the standard
axioms \texttt{propext}, \texttt{Classical.choice}, \texttt{Quot.sound}. PrintAxioms lines at each stage's final gate:

\begin{center}
\small
\begin{tabular}{lcccccc}
\toprule
\textbf{Stage} & \textbf{01} & \textbf{02} & \textbf{03} & \textbf{04} & \textbf{05} & \textbf{06}\\
\midrule
Standard / total & $161$ lines & $195/195$ & $519/519$ & $592/592$ & $747/747$ & $836/836$\\
\bottomrule
\end{tabular}
\end{center}

\begin{sloppypar}
Every deep input is a \textbf{registered hypothesis}: a named \texttt{Prop} passed as an explicit argument, never an
axiom. At \texttt{4405815} there are $29$ registered hypotheses ($28$ \texttt{Langlandscape.Hyp.*} rows plus the reused
\texttt{GL2TF.Hyp.EichlerMass}). Among them: \texttt{Hyp.SA}, \texttt{HS}, \texttt{BTtree}, \texttt{SerreCoh},
\texttt{SerreFree} (stage 2); \texttt{EllipticCount}, \texttt{IharaH0}
(stage 3); \texttt{HidaControl}, \texttt{HsiehCR}, \texttt{MordLambdaFree}, \texttt{InvLimExact}
(stage 4); \texttt{LambdaJL}; \texttt{Hyp.X.char}, \texttt{.iso}, \texttt{.gal} (stage 5); \texttt{HsiehGor}, \texttt{OldDecomp},
\texttt{HsiehGorFull}, \texttt{FreeRankTwo}, \texttt{TamSpec}, \texttt{HidaControlGL2}, \texttt{KimOta},
\texttt{KimOtaRT} (stage 6, open) and \texttt{ResidualBurnside} (conditionally discharged).
\textbf{Discharged} (3): \texttt{Hyp.PadicOrdDecomp}, \texttt{Hyp.RamNrdSurj} (\texttt{Langlandscape.Hyp.ramNrdSurj\_holds}, merged on main \texttt{2323a252}; drafted by an external model, run by Fabian; integrated, audited and gated by the agents) and \texttt{Hyp.BrandtStabFree} (\texttt{Langlandscape.Hyp.brandtStabFree\_holds}; drafted by an external model on Abhijit's machine; integrated, audited (Fable, PASS) and gated by the agents). Registered: $29$ (unchanged); conditionally discharged: $1$ (\texttt{ResidualBurnside}); open: $25$. \textbf{Final Lean gate}: Lean tree \texttt{c70ef3a9}, from scratch on TITAN (bsf-gate-1), PrintAxioms $955/955$ standard (previous gates: $950/950$ at \texttt{28d2b091}, $840/840$ at \texttt{2323a252}). Toy instances in the tests show consistency (non-vacuity) only.
\end{sloppypar}

\end{document}
